\rtmtight
\begin{tblr}{width=\textwidth,colspec={|Q[c,m,2.1cm]|X[l,m]|},hlines,vlines,hspan=minimal,rowsep=1.5pt,colsep=3pt,row{1}={bg=headgray,font=\bfseries}}
\SetCell{c,m}\hfil\logo\hfil & \textbf{POLITEKNIK NEGERI MALANG}\par\textbf{JURUSAN TEKNOLOGI INFORMASI}\par\textbf{PROGRAM STUDI: D4 TEKNIK INFORMATIKA} \\
\SetCell[c=2]{l} \textbf{RENCANA TUGAS MAHASISWA (RTM) DAN RUBRIK PENILAIAN --- DRAF KONSEP B: PENGABDIAN \& SIKAP} & \\
Nama Mata Kuliah & Kuliah Kerja Nyata Tematik \\
Kode Mata Kuliah & RTI257002 (sementara) \quad \textbf{sks / Jam perminggu:} 2 SKS / 4 jam/minggu \quad \textbf{Semester:} 6/7 (sementara) \\
Dosen Pengampu & Tim Pengajar Program Studi D4 TI \\
\end{tblr}
\assessmentblock{Proposal Program Kerja}{Case Method dan penyusunan dokumen}{SCPMK0101-06001, SCPMK0301-06003, SCPMK0302-06004}{Mahasiswa menyusun proposal program kerja pemberdayaan berdasarkan observasi dan analisis kebutuhan sosial mitra/komunitas, memuat rencana kegiatan, jadwal, dan target dampak sosial.}{Melakukan observasi dan sosialisasi lapangan, menganalisis kebutuhan sosial secara partisipatif, menyusun proposal, dan mempresentasikannya dalam seminar proposal.}{Dokumen proposal program kerja pemberdayaan dan bahan presentasi seminar proposal.}{Ketepatan analisis kebutuhan sosial dan kelayakan program & 7 & Analisis kebutuhan sosial akurat dan mendalam; program kerja yang diusulkan realistis, relevan, dan menjawab masalah nyata komunitas. & Analisis kebutuhan cukup memadai; program kerja relevan tetapi kurang detail pada beberapa aspek pelaksanaan. & Analisis kebutuhan dangkal atau tidak berbasis data lapangan; program kerja tidak relevan atau tidak realistis. \\ \\
Kejelasan target dampak dan rencana kerja & 7 & Target dampak sosial dirumuskan secara terukur, rencana kerja rinci dan layak dilaksanakan bersama mitra. & Target dampak cukup jelas tetapi beberapa bagian rencana kerja masih perlu disempurnakan. & Target dampak tidak jelas atau tidak terukur; rencana kerja tidak layak dilaksanakan. \\ \\
Kualitas presentasi dan etika penyusunan & 6 & Presentasi proposal terstruktur dan meyakinkan; etika penyusunan (kejujuran data, pengakuan sumber) terjaga penuh. & Presentasi cukup baik dengan sedikit kekurangan struktur; etika penyusunan secara umum terjaga. & Presentasi tidak terstruktur atau terdapat pelanggaran etika penyusunan proposal. \\}

\assessmentblock{Logbook dan Keterlibatan Sosial}{Portofolio dan praktik lapangan}{SCPMK0102-06002, SCPMK0302-06004}{Mahasiswa mengisi logbook kegiatan harian selama pelaksanaan program pemberdayaan yang mencatat aktivitas, keterlibatan, dan refleksi kontribusi sosial secara berkala.}{Mengisi logbook harian, mendokumentasikan kegiatan pemberdayaan, melakukan refleksi mingguan, dan mendapatkan validasi mitra/DPL.}{Logbook digital/fisik yang terisi lengkap, divalidasi mitra/DPL, dan dokumentasi kegiatan pemberdayaan.}{Kelengkapan dan konsistensi pengisian logbook & 10 & Logbook terisi setiap kegiatan, mencakup aktivitas, waktu, hasil, dan refleksi; divalidasi mitra/DPL secara berkala. & Logbook terisi sebagian besar kegiatan dengan informasi yang cukup, tetapi beberapa entri tidak lengkap atau tidak tervalidasi. & Logbook banyak kegiatan kosong, isian tidak informatif, atau tidak divalidasi mitra/DPL. \\ \\
Tingkat keterlibatan dan kontribusi nyata dalam pemberdayaan & 10 & Mahasiswa terlibat aktif dan konsisten dalam seluruh rangkaian kegiatan pemberdayaan, memberikan kontribusi nyata yang diapresiasi komunitas. & Keterlibatan cukup aktif tetapi tidak konsisten sepanjang program; kontribusi ada namun terbatas pada kegiatan sederhana. & Keterlibatan minim, mahasiswa pasif dalam kegiatan, atau tidak memberikan kontribusi nyata bagi komunitas. \\ \\
Kualitas refleksi tanggung jawab sosial & 10 & Refleksi menunjukkan pemahaman mendalam terhadap tanggung jawab sosial, menghubungkan pengalaman dengan nilai pengabdian masyarakat. & Refleksi cukup informatif tetapi hubungan dengan tanggung jawab sosial masih terbatas. & Refleksi bersifat deskriptif semata tanpa analisis atau koneksi dengan nilai pengabdian masyarakat. \\}

\assessmentblock{Penilaian Mitra/DPL}{Penilaian sikap dan dampak sosial}{SCPMK0101-06001, SCPMK0102-06002}{Penilaian dari mitra/komunitas dan dosen pembimbing lapangan (DPL) terhadap sikap, etika kerja, dan dampak sosial yang dihasilkan mahasiswa selama pelaksanaan KKN.}{Mitra/DPL mengobservasi sikap dan kinerja harian, mengevaluasi kontribusi terhadap komunitas, memberikan penilaian mid-term dan akhir menggunakan formulir penilaian standar.}{Formulir penilaian mitra/DPL mid-term dan akhir yang telah diisi dan ditandatangani.}{Profesionalisme, disiplin, dan etika kerja & 10 & Mahasiswa konsisten hadir tepat waktu, bersikap sopan dan menghargai norma komunitas, serta mematuhi kesepakatan program tanpa pengingat. & Sikap dan profesionalisme umumnya baik tetapi kadang ada keterlambatan atau pelanggaran kecil kesepakatan program. & Sering terlambat, tidak menghargai norma komunitas, atau menunjukkan sikap tidak profesional berulang kali. \\ \\
Dampak nyata program bagi mitra/komunitas & 8 & Program pemberdayaan memberikan dampak nyata dan diakui bermanfaat oleh mitra/komunitas, melebihi ekspektasi awal. & Program memberikan dampak yang cukup dirasakan komunitas meski belum optimal pada seluruh sasaran. & Program tidak memberikan dampak yang jelas atau tidak dirasakan manfaatnya oleh komunitas. \\ \\
Kemampuan beradaptasi dengan konteks sosial mitra & 7 & Mahasiswa mampu beradaptasi dengan cepat dan sensitif terhadap norma sosial-budaya mitra, membangun kepercayaan komunitas. & Adaptasi cukup baik meski memerlukan waktu; kepercayaan komunitas terbangun secara bertahap. & Mahasiswa kesulitan beradaptasi dengan konteks sosial mitra atau menimbulkan kesalahpahaman dengan komunitas. \\}

\assessmentblock{Laporan dan Seminar Hasil}{Case Method, penulisan ilmiah, dan presentasi}{SCPMK0301-06003, SCPMK0102-06002}{Mahasiswa menyusun laporan KKN yang komprehensif serta mempresentasikan dan mempertahankan hasil program pemberdayaan di hadapan tim penguji.}{Mengumpulkan data dan dokumentasi selama pelaksanaan program, menyusun laporan sesuai sistematika Polinema, bimbingan dengan DPL, revisi, dan presentasi seminar hasil.}{Naskah laporan KKN lengkap beserta lampiran (logbook, dokumentasi dampak sosial), dan bahan presentasi seminar hasil.}{Kelengkapan struktur dan dokumentasi dampak sosial & 9 & Laporan memuat semua bab yang dipersyaratkan dengan dokumentasi dampak sosial program yang lengkap dan akurat. & Struktur laporan lengkap tetapi beberapa bagian dokumentasi dampak sosial kurang mendalam. & Struktur laporan tidak lengkap atau dokumentasi dampak sosial tidak mencerminkan kegiatan yang sesungguhnya. \\ \\
Kualitas presentasi dan komunikasi hasil & 8 & Presentasi terstruktur, disampaikan dengan percaya diri, dan mampu menyampaikan dampak sosial program secara meyakinkan kepada penguji. & Presentasi cukup terstruktur dan dapat dipahami, tetapi penyampaian kurang percaya diri atau kurang meyakinkan. & Presentasi tidak terstruktur, bahasa tidak profesional, atau materi sulit dipahami penguji. \\ \\
Kemampuan menjawab pertanyaan dan refleksi kritis & 8 & Semua pertanyaan dijawab dengan tepat, didukung bukti dari pengalaman lapangan dan refleksi kritis atas peran mahasiswa dalam komunitas. & Sebagian besar pertanyaan dijawab dengan cukup baik, tetapi beberapa jawaban kurang didukung bukti atau refleksi. & Banyak pertanyaan tidak dapat dijawab, jawaban tidak relevan, atau tidak dapat mempertahankan konten laporan. \\}

\begin{tblr}{width=\textwidth,colspec={|Q[l,m,3.2cm]|X[l,m]|},hlines,vlines,hspan=minimal,row{1-3}={bg=headgray},rowsep=1.5pt,colsep=3pt}
Jadwal Pelaksanaan: & Minggu 1--16 sesuai RPS. Setiap evaluasi memiliki RTM dan rubrik multi-kriteria. \\
Lain-Lain Yang Diperlukan: & LMS, logbook digital, perangkat presentasi, dan alat peraga kegiatan komunitas. \\
Pustaka: & Mengacu pustaka utama dan pendukung pada bagian RPS. \\
\end{tblr}
